\begin{exsol}{14}{\MR{} Food for Thought 4.2}{119}{Exam}{A3}
\asked Identify the core competencies behind Swiss dominance of mechanical watchmaking, explain
how battery-powered electronic watches changed them, and say what else the new watchmakers
produce.

\method Core competencies as scale economies in development shared across products, and economies
of scope (\MR{} Ch.~4, pp.~107--108, with the Casio example on p.~107). Complementarities, coherent
systems and the shift of the best system (pp.~108--113, Figure~4.3). Innovation attributes (p.~92,
pp.~112--113). Ford's coherent but superseded strategy as the parallel case (p.~109, p.~111).

\solution
\subparagraph{The Swiss competencies.} Accuracy in a spring-driven watch is a mechanical
achievement. The competencies were (i) the design of micro-mechanical movements; (ii) precision
fabrication and finishing of tiny parts (springs, gears, escapements, jewelled bearings);
(iii) highly skilled hand assembly and regulation, reproduced by long apprenticeship; (iv) a dense
network of specialized small suppliers coordinated by long-standing relationships; and (v) a
reputation for quality and a distribution channel through jewellers. These pieces formed a
\emph{coherent system} (pp.~111--112): small batches, craft labour paid for skill, quality
achieved by the worker rather than by the machine, high prices justified by brand and
craftsmanship. Each element raised the return to the others --- complementarity in the sense of
p.~108.

\subparagraph{What the new technology changed.} With a quartz oscillator driven by an integrated
circuit and a battery, accuracy comes from electronics, and it is better and cheaper than the best
hand-regulated movement. The relevant competencies became semiconductor design and fabrication,
displays (LCD), miniature batteries, and high-volume automated assembly. These have very large
scale economies, and they are \emph{shared inputs} across many products: exactly the core
competencies of p.~107. Casio's case in the text illustrates the logic: scale economies in LCDs
made calculators, wristwatches and electronic address books cheaper together, and large calculator
sales lowered display costs and made watch entry profitable (p.~107). The value of the Swiss
competencies for the mass market fell sharply: mechanical skill no longer bought accuracy, and the
craft network could not match semiconductor-scale costs.

\subparagraph{Why the incumbents struggled to adapt.} In Figure~4.3's terms (pp.~111--113),
technology moved the profit landscape so that a different coherent system --- mass production of
electronic modules, automated assembly, low prices, wide distribution --- became the best one.
Managers inside the craft system can adjust \emph{locally} (cheaper springs, better tooling), but
local moves keep them on the old peak. Reaching the new one requires changing many complementary
elements at once --- skills, equipment, suppliers, distribution, pricing --- and partial moves
lose money, the ``valley'' of the Aula~1 notes. It also requires information no one in the
organization has (innovation attributes, pp.~92, 113): watchmakers did not know semiconductors. In
short, the competencies that made the old system excellent were not only useless for the new one,
they were \emph{complementary to the old system}, which made the jump harder.

The old system did not become worthless everywhere: it remained coherent for a luxury segment
where a mechanical watch is valued as craftsmanship and jewellery rather than as a timekeeper.
Broadly, the Swiss industry's later recovery paired a reorganized, low-cost quartz line with this
repositioning of mechanical watches --- two coherent systems for two segments.

\subparagraph{Other products of the new watchmakers.} Firms whose competencies are in
electronics apply them across product lines: calculators, electronic organizers and address books
(p.~107), and more generally consumer electronics, displays and the semiconductor components
themselves. For them the watch is one more product that shares components and development
capabilities with the rest --- an economy of scope, not a craft.

\interp A coherent system is a strength only for the environment it was built for (p.~109: Ford's
strategy was ``equally coherent'' for its time). New technology can make a different coherent
system optimal, and complementarities turn yesterday's competencies into today's inertia. Choosing
the new pattern is a central, strategic decision with innovation attributes (p.~115).

\verify Alternative argument by comparison: the same mechanism explains Ford's loss to GM and
GM's to Toyota (pp.~16, 111). In each case the incumbent's system stayed internally consistent,
while the environment changed which system was best. Limit case: an innovation that improved the
spring or the escapement would have \emph{raised} the value of Swiss skills and called only for
local adjustment. What made the quartz watch destructive is that it moved the source of accuracy
outside the competencies the system was built on.

\answer{Swiss competencies were micro-mechanical design, precision part-making, skilled hand
assembly, a network of specialized suppliers and a quality brand --- a coherent craft system. The
quartz-electronic watch moved accuracy into semiconductors, displays and batteries made at large
scale, the core competencies of electronics firms, which also make calculators, organizers and
other consumer electronics. Complementarities and missing information made the Swiss jump to the
new system slow; the craft system survived mainly in the luxury segment.}
\end{exsol}
